\documentclass{situationist-preprint}

\title{Non-Deviant Causation}
\msid{ACT-03-a}
\msfamily{ACT-03}
\msdiscipline{Action, Agency and Free Will}
\mstype{Principal manuscript}
\mskeywords{action; causal theory of action; deviant causal chains; Davidson; guidance; counterfactual sensitivity; control}

\begin{document}
\maketitle

\begin{abstract}
What distinguishes an action from a mere bodily movement? The causal theory answers that an action is a movement caused by a reason, and Davidson's climber, who is so unnerved by his wish to be rid of his companion that his grip loosens, shows that the cause must operate \enquote{in the right way}. Davidson doubted that the right way could be specified. We argue that it can: an action is a bodily movement caused by a reason and sensitive to counterfactual variations in that reason, in the sense that, over a range of nearby alternative contents, the movement would have varied as the alternative reason prescribes, holding fixed whatever is idle in the actual sequence. We show that sensitivity entails counterfactual dependence and is strictly stronger, that no causal route through a content-insensitive state such as agitation can be sensitive, and that any intermediary preserving the reason's rational distinctions can. The classical cases of basic and consequential deviance fail the condition; Frankfurt interveners, prosthetic relays and ideomotor action satisfy it. The condition specifies guidance without using the notion of action, and it fits the reasons-responsiveness account of free will defended in ACT-01: sensitivity is required of every action, responsiveness of every free one.
\end{abstract}

\section{The problem}\label{s:problem}

\subsection{The question}

When I raise my arm, my arm rises. Wittgenstein asked what is left over if we subtract the fact that my arm goes up from the fact that I raise it \citep[\S 621]{Wittgenstein1953}. Two answers have dominated. The causal theory says that what is left over is a cause: the movement is an action because it is caused by a reason of the agent's, a desire for some end together with a belief that the movement is a way of attaining it \citep{Davidson1963}. Its rivals say that what is left over is not a cause but something else: the intelligibility of the movement under a description \citep{Anscombe1957,Melden1961}, its teleological structure \citep{Wilson1989,Sehon2005}, or the agent's own causing of it \citep{Hyman2015,Steward2012}.

The causal theory has a well-known difficulty. A reason may cause a movement that it rationalizes without the movement's being an action. Davidson's climber wants to rid himself of the weight and danger of holding another man on a rope, and believes that loosening his grip would do so; this belief and desire so unnerve him that his grip loosens \citep{Davidson1973}. The reason caused the movement, and the movement is the one the reason rationalizes, but the climber did not loosen his grip; it loosened. The cause must operate in the right way, and Davidson despaired of saying what the right way is without using the notion of acting intentionally, which was to be explained.

We argue that the right way can be specified. A reason causes a movement in the right way when the movement is \emph{sensitive to counterfactual variations in the reason}: had the reason's content been different in any of a range of nearby respects, the movement would have been different, and different in the way the varied reason prescribes. The climber's movement is not sensitive in this sense. Any vivid thought of what to do with the rope would have unnerved him, and an unnerved climber loosens his grip whatever he was thinking about doing with it.

\subsection{A short history}\label{s:history}

\paragraph{The causal theory and its critics.} \citet{Davidson1963} defended two theses against the post-Wittgensteinian consensus represented by \citet{Melden1961}: that reasons are causes, and that rationalizing explanation is a species of causal explanation. The argument for the second thesis was that an agent may have a reason to do something and do it, and yet not do it for that reason; what makes the reason the one she acted on is that it caused the action. \citet{Goldman1970} developed the theory in detail, and by the 1970s it was the standard view. Its critics did not deny that reasons are causes so much as deny that causation could do the work of the \enquote{because} in \enquote{she did it because she wanted to}. The deviant cases became their main weapon.

\paragraph{Deviance.} Three kinds of case were distinguished. In \emph{basic} deviance, the chain from reason to bodily movement is deviant, as with the climber. \citet{Frankfurt1978} gave another: a man at a party intends to spill his drink as a signal, and his intention so agitates him that his hand trembles and the drink spills. In \emph{consequential} deviance, the movement is an action but the chain from the movement to its intended result is deviant: a man shoots to kill, misses, and the shot stampedes a herd of wild pigs that tramples the victim. \citet{Chisholm1966} described a nephew whose intention to kill his uncle so agitates him that he drives recklessly and kills a pedestrian, who turns out to be the uncle. In \emph{antecedential} deviance, the chain within the agent's reasoning, before any movement, is deviant: a reason causes an intention by a route that bypasses practical reasoning.

\paragraph{Proximal intentions.} The first response was to require that the reason cause the movement through an intention, and in particular through a proximal intention or \enquote{intention in action} that is present while the movement occurs. \citet{Searle1983} added that the intention in action is causally self-referential: its content is that this very intention cause the movement. \citet{Brand1984} and \citet{Mele1992} developed related views on which a proximal intention guides the movement as it unfolds. The response handles cases in which the deviance occurs between the reason and the formation of the intention. It does not handle Frankfurt's party guest, whose proximal intention is what agitates him, or the climber if he has already decided to let go and the decision itself unnerves him. Self-reference does not help either: the content of an intention that it cause the movement is satisfied by a deviant chain, since the intention does cause it.

\paragraph{Guidance.} \citet{Frankfurt1978} proposed that the distinguishing mark of action is not the movement's causal history but the agent's guidance of it while it occurs: a movement is an action if it is under the guidance of the agent, in the sense that mechanisms are prepared to compensate for perturbations that would divert it from its course. The proposal is the right one in spirit, and our thesis is a precise version of it. But Frankfurt stated guidance in terms of compensation by the agent, which threatens to presuppose agency, and he offered the account as an alternative to the causal theory rather than a completion of it.

\paragraph{Differential explanation.} \citet{Peacocke1979} gave the first precise proposal of the kind we defend: the intention must \emph{differentially explain} the movement, in that there is a law-like dependence of the movement's parameters on the parameters of the intention's content, so that small variations in what is intended would be matched by variations in what is done. \citet{Bishop1989} added that the dependence must be sustained by feedback through the agent's own perceptual and motor systems, so as to exclude relays through another agent's decisions. The differential explanation view was criticized on two grounds. First, it seemed to exclude actions whose parameters are coarse, such as pressing a button or uttering a word, which have no continuous variation. Second, deviant chains could be constructed in which variations happened to be matched, so that the movement covaried with the content by luck or by a wayward but systematic route. We answer both criticisms below (Sections~\ref{s:cases} and~\ref{s:objections}).

\paragraph{Reliability and repertoire.} More recent work has tied the right way to reliability. \citet{Schlosser2007} argued that in basic deviance the mental state does not cause the movement in virtue of its content, so that the content is not causally relevant to what happens; the right way is causation in virtue of content. \citet{Aguilar2012} proposed that a movement is an action when it is produced by a reliable exercise of the agent's action repertoire, the set of movements she can bring about reliably by intending them. \citet{Shepherd2021} made control the central notion of agency and analysed it as the reliable and flexible success of the agent's plans across a range of circumstances. Our condition is the common core of these proposals. Causation in virtue of content is, on any interventionist account of causal relevance \citep{Woodward2003}, a matter of how the effect would vary under variations in the content; reliability over a repertoire is variation that matches the content across the repertoire; and flexibility across circumstances is the same matching with different variables held fixed.

\paragraph{The non-causal alternatives.} The Anscombean tradition holds that what makes a movement an action is that the question \enquote{Why?} in a certain sense has application to it, and that the answer gives a description under which it is intentional \citep{Anscombe1957,Thompson2008}. Teleological theorists hold that action explanation is irreducibly teleological \citep{Wilson1989,Sehon2005}. Agent-causal theorists hold that an action is the agent's causing of a movement, and that the agent is not identical with any of her states \citep{Hyman2015}. We do not argue against these views directly. We argue that the deviant cases, which are their main argument against the causal theory, can be answered. If they can, the remaining disagreement concerns whether the causal theory captures the agent's own role in action \citep{Velleman1992}, and Section~\ref{s:objections} takes up that challenge.

\subsection{Plan}

Section~\ref{s:framework} states the framework, Section~\ref{s:argument} the main argument and three structural results. Section~\ref{s:cases} applies the account to the standard cases and to some new ones. Section~\ref{s:consequences} relates it to the theories of free will, akrasia and mental causation defended elsewhere in this series. Sections~\ref{s:objections} and~\ref{s:limits} give objections and limits.

\section{Framework}\label{s:framework}

We work with causal models of the kind used in MET-04-a: variables, structural equations, and interventions that set some variables and leave the rest to their equations. MET-04 argues that a cause is a counterfactual difference-maker selected by norms of normality, and we use that notion of cause. The one novelty is that some variables take as values the \emph{contents} of the agent's mental states.

\begin{definition}[Reason]\label{d:reason}
A \emph{reason} of an agent, in the motivating sense, is a pair of states: a pro-attitude toward an end and a belief that a bodily movement of a certain type is a way of attaining it, in the agent's circumstances. The \emph{content} of the reason is the pair of contents of these states. A \emph{reason variable} is a variable whose values are possible contents of a reason the agent might have in the situation.
\end{definition}

\begin{definition}[Rationalized movement]\label{d:rat}
For a content $c$ of a reason, $\rho(c)$ is the type of bodily movement that the reason with content $c$ represents as the way to its end. A reason with content $c$ \emph{rationalizes} a movement of type $\rho(c)$.
\end{definition}

The map $\rho$ is fixed by the contents, not by the agent's causal constitution: it records what the agent's reason calls for, whether or not her body would oblige.

\begin{definition}[Variation range]\label{d:range}
A \emph{variation range} for a reason with actual content $c$ is a set $V$ of contents that differ from $c$ in respects the agent is capable of acting on (the target of a reach, the force of a grip, the choice between pressing one button and another), and that she would have in situations close to the actual one. The range is \emph{non-trivial} if it contains a content $c'$ with $\rho(c') \neq \rho(c)$.
\end{definition}

\begin{definition}[Profile]\label{d:profile}
Let $H$ be the set of variables that are idle in the actual sequence from the reason to the movement: variables, such as the trigger of a counterfactual intervener, whose actual values play no role in producing the movement. The agent's \emph{profile} is the function $M$ that assigns to each content $c'$ the type of movement that occurs under the intervention that sets the reason variable to $c'$ and every variable in $H$ to its actual value.
\end{definition}

\begin{definition}[Sensitivity]\label{d:sens}
A movement is \emph{sensitive to counterfactual variations in a reason} with actual content $c$ if there is a non-trivial variation range $V$ for the reason such that $M(c') = \rho(c')$ for every $c'$ in $V \cup \{c\}$.
\end{definition}

\begin{definition}[Action]\label{d:action}
A bodily movement is an \emph{action} if and only if (i) it is caused by a reason of the agent's, in the sense of MET-04, and (ii) it is sensitive to counterfactual variations in that reason.
\end{definition}

Definition~\ref{d:action} is the thesis. The argument for it rests on two principles.

\begin{principle}[Control]\label{p:control}
A bodily movement is an action if and only if it is caused by a reason of the agent's and is under that reason's control.
\end{principle}

\begin{principle}[Control as covariation]\label{p:cov}
A state with content controls an outcome if and only if, holding fixed what is idle in the actual sequence, the outcome would vary with the state's content over a range of nearby contents, in the way that the content prescribes.
\end{principle}

Principle~\ref{p:control} is the common ground of the causal theory and of Frankfurt's guidance view: the climber's grip is not under the control of his reason, although it is caused by it, and that is why it is not an action. It is what remains of the causal theory once the deviant cases are granted. Principle~\ref{p:cov} is the substantive claim. It is the ordinary engineering notion of control: a thermostat controls the room's temperature because the temperature would track the setting over a range of settings, and a pilot controls an aircraft because its heading would follow the yoke over a range of positions. Mere dependence is not control: a room whose heating comes on whenever the thermostat is touched, at whatever setting, depends on the thermostat but is not controlled by it.

\begin{assumption}[Contentful causes]\label{a:content}
Reason variables can be causes in apt causal models: interventions on the content of a reason are well defined, and the movement's dependence on them is a matter of fact.
\end{assumption}

Assumption~\ref{a:content} is what the causal theory needs from the philosophy of mind. It is supported by MND-11, which argues that mental causes are proportional causes, so that a model whose variables are contents, rather than neural states, is the apt one when the question is what an agent's movement depends on.

\section{The argument}\label{s:argument}

\begin{theorem}\label{t:main}
A bodily movement is an action if and only if it is caused by a reason and sensitive to counterfactual variations in that reason.
\end{theorem}

\begin{proof}
By Principle~\ref{p:control}, a movement is an action if and only if it is caused by a reason and under that reason's control. By Principle~\ref{p:cov}, a reason controls a movement if and only if, holding fixed what is idle in the actual sequence, the movement would vary with the reason's content over a range of nearby contents in the way the content prescribes. A content prescribes the movement that it rationalizes (Definition~\ref{d:rat}); nearby contents form a variation range (Definition~\ref{d:range}); the movement's variation, holding fixed what is idle, is given by the profile (Definition~\ref{d:profile}). Since a range on which the movement does not vary at all is not one over which it is controlled, the range must be non-trivial. Hence the reason controls the movement if and only if the movement is sensitive to it (Definition~\ref{d:sens}).
\end{proof}

The theorem would be of little interest if sensitivity were simply dependence. It is not. The following results show that it is strictly stronger, and identify what it excludes.

\begin{lemma}\label{l:dep}
If a movement is sensitive to counterfactual variations in a reason, it counterfactually depends on the reason's content, holding fixed the variables in $H$.
\end{lemma}

\begin{proof}
Let $V$ be a non-trivial range witnessing sensitivity, and $c' \in V$ with $\rho(c') \neq \rho(c)$. Then $M(c') = \rho(c') \neq \rho(c) = M(c)$, so setting the reason variable to $c'$, with $H$ at its actual values, yields a different movement.
\end{proof}

Lemma~\ref{l:dep} shows that clause (ii) of Definition~\ref{d:action} entails de facto difference-making in the sense of MET-04-a; clause (i) adds the normality condition that MET-04 places on causes. The converse of the lemma fails, and the climber is the counterexample.

\begin{proposition}\label{p:constant}
If the profile is constant on $V \cup \{c\}$ for a non-trivial range $V$, the movement is not sensitive to the reason over $V$. In particular, a movement can depend counterfactually on the presence of a reason without being sensitive to it.
\end{proposition}

\begin{proof}
If $M$ is constant on $V \cup \{c\}$ and the movement were sensitive over $V$, then for the $c' \in V$ with $\rho(c') \neq \rho(c)$ we would have $\rho(c') = M(c') = M(c) = \rho(c)$, a contradiction. For the second part, let the reason variable have, besides contents, a value $\bot$ for the absence of any relevant reason, and let $M(\bot)$ be no movement and $M(c') = \rho(c)$ for every content $c'$. The movement depends on the presence of the reason, since $M(\bot) \neq M(c)$, but it is not sensitive to it over any non-trivial range.
\end{proof}

The second model is the climber. Had he had no thought of the rope, he would not have been unnerved, and his grip would not have loosened: the loosening depends on his reason. But had his reason been to secure his grip, or to shift the rope to the other hand, the thought would have unnerved him equally, and his grip would have loosened just the same. His profile is constant over the contents that matter.

Many deviant chains have a more specific structure: the reason affects the movement only through an intermediate state.

\begin{proposition}\label{p:inter}
Suppose the reason affects the movement only through an intermediate variable $I$, so that $M(c') = h(g(c'))$, where $g$ gives the value of $I$ under each content and $h$ the movement under each value of $I$. Then the movement is sensitive over $V$ if and only if $h \circ g$ agrees with $\rho$ on $V \cup \{c\}$. It follows that sensitivity requires $g$ to separate every two contents in $V \cup \{c\}$ that rationalize different movements, and that any intermediary satisfying this requirement, combined with a suitable $h$, yields sensitivity.
\end{proposition}

\begin{proof}
The first claim restates Definition~\ref{d:sens} with $M = h \circ g$. If $g(c_1) = g(c_2)$ while $\rho(c_1) \neq \rho(c_2)$, then $M(c_1) = M(c_2)$, so $M$ cannot agree with $\rho$ on both. Conversely, if $g$ separates all such pairs, define $h$ on the image of $g$ by $h(g(c')) = \rho(c')$; this is well defined because contents with the same image rationalize the same movement, and $h \circ g$ then agrees with $\rho$.
\end{proof}

\begin{corollary}\label{c:agitation}
A causal chain from reason to movement that passes only through a state whose value does not distinguish between contents rationalizing different movements, such as agitation, nervousness or a tremor, does not make the movement sensitive to the reason, and so does not make it an action.
\end{corollary}

\begin{proof}
Agitation is a state whose value (its presence and intensity) is fixed by how vivid or distressing the reason is, not by which movement it rationalizes. Two contents of equal vividness that call for different movements yield the same value of $I$. By Proposition~\ref{p:inter}, the movement is not sensitive over any range containing them, and the actual range does contain such contents.
\end{proof}

Proposition~\ref{p:inter} has a converse side that matters as much as the corollary. It does not care what the intermediary is. Any state that preserves the rational distinctions among the contents in the range, and any process that maps it onto the movements those contents call for, carries control. The account therefore puts no restriction on the physical realization of the route from reason to movement: a damaged motor pathway replaced by a prosthesis, or an unusual route through the spinal cord, serves as well as the normal one. This is the point on which our account differs from Bishop's requirement of feedback through the agent's own systems, and we defend it in Section~\ref{s:cases}.

\begin{remark}\label{r:coarse}
Definition~\ref{d:sens} does not require that the movement have continuous parameters. A variation range may consist of finitely many contents, such as the intention to press the left button and the intention to press the right one. What is required is that the movement would have matched each of them. This answers the first criticism of \citet{Peacocke1979}: coarse actions are sensitive to coarse variations.
\end{remark}

\section{Cases}\label{s:cases}

\paragraph{Basic deviance.} The climber and Frankfurt's party guest are handled by Corollary~\ref{c:agitation}. In each case the route from reason to movement passes through agitation, and agitation does not distinguish between the reasons that call for different movements. Had the guest intended to signal by raising his glass rather than spilling it, his hand would have trembled and the drink spilled all the same.

\paragraph{Consequential deviance.} Definition~\ref{d:action} concerns bodily movements, but the same condition, applied to the result of a movement, distinguishes results the agent brings about intentionally from results she merely causes.

\begin{remark}\label{r:result}
A result of an action is brought about intentionally if it is caused by the reason, through the action, and sensitive to counterfactual variations in the reason over a range of contents concerning that result.
\end{remark}

The killer's shot is an action: had he aimed elsewhere, the gun would have pointed elsewhere. The victim's death is not intentionally brought about. Had he aimed a little to the left or to the right, at a victim standing a little to the left or to the right, the pigs would have stampeded in the same direction, or not at all, depending on where they happened to be; the death does not vary with the aim. Chisholm's nephew is a case of basic deviance with a consequential sequel: his reckless driving is caused through agitation, and the death of the pedestrian, who happens to be his uncle, would not have followed a change in whom he intended to kill.

\paragraph{Antecedential deviance.} The present account does not concern the route from reasons to intentions. A reason may cause an intention by a route that bypasses deliberation, and the intention may then produce a sensitive movement; the movement is an action, though perhaps not one done for that reason. The question whether the agent acted for the reason that caused her intention, or for the intention alone, is one about the explanation of an action, not about whether something is one. Section~\ref{s:limits} records this limit.

\paragraph{Frankfurt interveners.} Black stands ready to make Jones raise his arm should Jones show signs of not raising it; Jones raises it on his own \citep{Frankfurt1969}. Jones's movement is an action, and ACT-02 argues that he is responsible for it. Naive counterfactual covariation gives the wrong answer: had Jones's reason been different, Black would have intervened and the arm would have risen anyway. Definition~\ref{d:profile} prevents this. Black's trigger is idle in the actual sequence, so it is held fixed at its actual value, which is that he does not intervene; with the trigger so held, the arm would have followed the reason. This is the same move as in MET-04-a's treatment of preemption, and it is warranted for the same reason: the counterfactuals that measure control are those that concern the actual route, not those that concern backup routes.

\paragraph{Relays and prostheses.} A patient whose motor nerve is severed has a device that reads her motor cortex and stimulates the muscles of her arm accordingly. Her reaching is an action, and by Proposition~\ref{p:inter} it satisfies Definition~\ref{d:action}, provided the device separates the contents that call for different reaches. Now replace the device with a neuroscientist who reads her intentions from a monitor and presses buttons that move the arm correspondingly, as reliably as the device. Cases of this form were discussed by \citet{Peacocke1979} and \citet{Bishop1989}, and they divide intuitions. Our verdict is that the patient's movement is her action, provided the neuroscientist's relaying is as reliable over the range as the device. If he relays only because he happens to be in a cooperative mood, and would have done otherwise had his mood been different, then his mood is not idle in the actual sequence, the profile must allow it to vary, and over the nearby range in which his mood differs, the arm does not follow the patient's reason. Whether a relay through another agent preserves control depends on whether that agent's cooperation is robust, not on whether the relay is an agent. This seems to us the correct verdict and the source of the divided intuitions: one imagines the cooperative neuroscientist either as a fixed part of the apparatus, in which case the patient acts, or as a free agent who might at any moment decline, in which case the patient merely causes him to act.

\paragraph{Ideomotor action.} \citet[ch.~26]{James1890} described movements that follow immediately on the thought of them, without any separate act of will: one thinks of the pen on the table, and one's hand is reaching for it. Such movements are actions on our account, since they vary with the content of the thought: one would have reached for the cup had one thought of the cup. This is the right verdict. It shows that the content-sensitive route need not include a deliberate decision, and that what makes the climber's case deviant is not the absence of a decision but the content-insensitivity of the route.

\paragraph{Skilled and habitual action.} An experienced driver brakes when the car ahead slows, without deliberating. The movement is an action, and it is sensitive: had she wanted to stop at the next junction, or to change lanes, she would have braked or steered accordingly. Habit makes the reason's operation automatic without making the movement insensitive to it. A tic, by contrast, may be triggered by a desire to suppress it, and has the profile of the climber.

\paragraph{Wayward covariation.} Suppose the climber is constituted so that a vivid thought of loosening his grip makes it loosen, a vivid thought of tightening it makes it tighten, and a vivid thought of shifting the rope makes his hand shift it, all through an anxious reflex that bypasses his ordinary motor control. His profile matches $\rho$ over the range. On our account his movements are actions. Section~\ref{s:objections} defends this verdict.

\section{Recent literature}\label{s:recent}

\paragraph{Causal relevance of content.} \citet{Schlosser2007} argues that what goes wrong in basic deviance is that the mental state does not cause the movement in virtue of its content. We agree, and Proposition~\ref{p:inter} gives the claim a precise form: a state causes a movement in virtue of its content, in the sense relevant to action, when the route from it to the movement preserves the distinctions between contents that call for different movements. Without such a specification, the appeal to content-relevance is open to the reply that the climber's reason does cause the loosening in virtue of its content, since a reason with a less distressing content would not have unnerved him. Our condition distinguishes the two kinds of content-dependence: dependence on how distressing a content is, and dependence on which movement it calls for. Only the second is control.

\paragraph{Repertoires and reliability.} \citet{Aguilar2012} requires that the movement be produced by a reliable exercise of the agent's action repertoire. A repertoire is a set of movement types the agent can perform reliably by intending to; our variation range is the part of the repertoire that is close to the actual reason. Aguilar's condition and ours agree on the classical cases. They differ on single exercises of an unreliable capacity. An agent whose reaching is accurate only half the time, because of a neurological condition, and who reaches successfully on this occasion, has acted if her movement, on this occasion, would have followed her reason over the nearby range; reliability across occasions is evidence of this, but not a condition of it. Our condition concerns the counterfactual structure of the actual occasion, as Lemma~\ref{l:dep} requires.

\paragraph{Control and skill.} \citet{Shepherd2021} analyses control as the disposition to succeed reliably and flexibly at what one intends across a range of circumstances, and makes it the basis of a gradualist account of agency, skill and knowledge of action. Our Principle~\ref{p:cov} is the special case of his notion in which what varies is the content of the reason and what is held fixed is the circumstance. Varying the circumstance instead, while holding the reason fixed, gives the compensatory guidance emphasized by \citet{Frankfurt1978}: the arm would reach the cup whether the cup were here or there, because the agent corrects her movement as she goes. Both are dimensions of control, and an action will usually exhibit both; but only the first is required to distinguish action from mere movement, since a ballistic movement that, once launched, cannot be corrected is still an action if it was launched by a reason it would have followed.

\paragraph{The agent's own role.} \citet{Velleman1992} objected that the standard causal theory describes a sequence in which the agent plays no part: states and events cause movements, and the agent is a spectator. \citet{Hyman2015} and \citet{Steward2012} conclude that action involves causation by the agent herself. Our account does not add an agent cause. It locates the agent's role in the profile: what the agent contributes to an action is that her body would follow her reasons over a range. This is a property of the agent, not of any single state, since the profile is fixed by her whole constitution, and MET-11's defence of irreducible dispositions secures its reality. Whether this is enough to answer Velleman is taken up in Objection~\ref{o:agent}.

\section{Consequences}\label{s:consequences}

\paragraph{Free will.} ACT-01-a analyses free will as the capacity to act from a mechanism of practical reasoning that is reasons-responsive and one's own, where responsiveness concerns what the mechanism would do were the reasons available to the agent different. The two counterfactual conditions are distinct, and both are needed. Sensitivity, the condition of Theorem~\ref{t:main}, concerns what the movement would do were the agent's reason different; it relates a given reason to its execution. Responsiveness concerns what reasons the agent would act on were the reasons available to her different; it relates the world to the agent's choice of reason. A compulsive hand-washer's movements are sensitive to his desire to wash: he would have washed the other hand, or turned the tap the other way, had he wanted to. They are actions. But the mechanism that issues in the desire is not responsive to reasons against washing, and so his actions are not free. Conversely, a free agent whose arm has been taken over by a deviant chain does not act freely in loosening her grip, since she does not act at all. Every free action is an action, and not every action is free; the two conditions are, in this sense, independent levels of a single counterfactual conception of agency.

\paragraph{Akrasia.} ACT-04 argues that an agent can intentionally act against her better judgment. On the present account, the akratic agent acts on the reason supplied by the desire she judges inferior, and her movement is sensitive to that reason: had she wanted the cake on the left rather than the one on the right, she would have reached for it. That her better judgment is not the reason that controls her movement does not make the movement any less an action. This agrees with ACT-04-a's claim that akratic action is intentional and is done for a reason the agent recognizes.

\paragraph{Mental causation.} MND-11 argues that mental causes are proportional causes and physical causes overly specific. Theorem~\ref{t:main} gives a reason to expect this in the case of action. The variations that matter for action are variations in content, and many neural realizations correspond to one content; a model whose variables are neural states would find the movement depending on fine details that make no difference to which movement occurs, and would miss the systematic covariation with content that constitutes control.

\section{Objections and replies}\label{s:objections}

\begin{objection}[Wayward covariation]\label{o:wayward}
The reflexively anxious climber of Section~\ref{s:cases} satisfies Definition~\ref{d:action}, but his movements are paradigms of deviance: they happen to him, through anxiety, and he does not perform them.
\end{objection}

\begin{reply}
We accept the verdict and deny that it is counterintuitive once the case is fully described. The original climber is deviant because anxiety is a blunt instrument: it makes him do one thing whatever he was thinking of doing. The reconstituted climber's anxiety is not blunt. It is a mechanism that, over the whole range of things he might think of doing with the rope, does that thing. Such a mechanism is a way of executing reasons, and in describing it we have described a peculiar motor system, whose operation resembles James's ideomotor action rather than a tremor. If the reflex is narrow, matching only a few contents and producing a generic loosening otherwise, the range is narrow too, and the case becomes borderline; but borderline cases of action are to be expected, since control comes in degrees of extent. The intuition that the reconstituted climber does not act survives, we think, only by importing the bluntness of the original case.
\end{reply}

\begin{objection}[Circularity]\label{o:circ}
The definition uses the notion of a variation range, which requires contents the agent is \enquote{capable of acting on}. Capability to act presupposes the notion of action, and the specification of the right way is circular after all, as Davidson feared.
\end{objection}

\begin{reply}
The phrase in Definition~\ref{d:range} can be replaced by one that does not mention action: the range consists of contents whose rationalized movements are movements the agent's body can make. Nothing in the definition then uses the notion of acting. Davidson's worry was specifically that the right way could be specified only as \enquote{the way that makes it an action}. The present specification uses contents, movements, counterfactual dependence and the map $\rho$ from contents to the movements they call for. The last is fixed by the contents of the agent's beliefs and does not presuppose that she acts.
\end{reply}

\begin{objection}[Arational actions]\label{o:arational}
\citet{Hursthouse1991} described actions done from emotion but for no reason: rumpling a loved one's hair, kicking a door in rage. These are actions, but they are not caused by a belief--desire pair that represents the movement as a means to an end. The thesis that an action is a movement caused by a reason is false.
\end{objection}

\begin{reply}
Hursthouse's point is that such actions are not explained by a belief about means and a desire for a further end. It is not that they lack a pro-attitude and a belief altogether. The angry agent wants to strike at the door, which he takes to be the thing frustrating him, and he believes that kicking is a way of striking it. The pro-attitude in Definition~\ref{d:reason} may be toward the movement itself, under a description, and the belief may be the trivial one that this movement is a way of performing a movement so described. With reasons so understood, the kick is sensitive: had he been angry at the window, he would have kicked the window, or struck it differently. What distinguishes arational actions is not the absence of a reason in the motivating sense but the absence of a further end, and Definition~\ref{d:action} does not require one.
\end{reply}

\begin{objection}[The disappearing agent]\label{o:agent}
\citet{Velleman1992}: on any causal theory, an action is a sequence in which mental states cause movements. The agent does nothing over and above what her states do. Adding counterfactual conditions on the states does not bring her back.
\end{objection}

\begin{reply}
The objection assumes that the agent's contribution, if there is one, must be an item in the actual sequence. It need not be. A profile is a property of the agent: it is fixed by her whole psychological and motor constitution, and it is not identical with the property of any single state in the sequence. When an action occurs, the agent's reasons cause her movement and her constitution makes the movement follow her reasons. This is the agent's contribution, and it is what is missing when her body is moved by a deviant chain, by another agent who does not reliably relay her will, or by a reflex. Velleman's own remedy, which identifies the agent with a desire to act in accordance with reasons, is compatible with ours; but it is not needed to distinguish actions from mere movements, which is the question of this paper.
\end{reply}

\begin{objection}[Range indeterminacy]\label{o:range}
How wide must the variation range be, and how close must its contents be to the actual content? The definition provides no answer, and different choices give different verdicts.
\end{objection}

\begin{reply}
The range is fixed by the situation: it consists of the contents the agent would have had in situations close to the actual one (Definition~\ref{d:range}). Where these are few, the condition is weaker, and where many, stronger. Different choices of model may give different verdicts at the margins, and we accept that \enquote{action} is a vague predicate there, with borderline cases between the reflexive and the controlled. MET-18 argues that vagueness lies in language and knowledge, not in the world, and on that view there is a fact at the margin, though perhaps not a knowable one. The classical deviant cases are not marginal: their profiles are constant over every reasonable range.
\end{reply}

\section{What the result does not establish}\label{s:limits}

First, Theorem~\ref{t:main} concerns bodily actions. Mental actions, such as calculating in one's head or deciding, and omissions, which ACT-07 treats, are not bodily movements, and whether a counterpart of the condition holds for them is not settled here. Second, the theorem does not settle antecedential deviance: it says when a movement is an action, not for which reason it was done when several reasons were present. Third, the theorem does not establish the causal theory against agent-causal rivals; it removes their principal argument, the deviant cases, but an agent-causalist may accept every result of Section~\ref{s:argument} and hold that the profile is an agent's power. Fourth, the theorem says nothing about whether the actions it identifies are free or responsible; those questions belong to ACT-01 and ACT-02. Finally, the results depend on Assumption~\ref{a:content}: if interventions on contents are not well defined, the counterfactuals in Definition~\ref{d:profile} have no determinate truth values, and the account must be restated for the neural realizers of contents, with the loss of proportionality noted in Section~\ref{s:consequences}.

\section{Conclusion}

A movement caused by a reason is an action when it is under the reason's control, and control is counterfactual covariation in the way the reason prescribes, holding fixed what is idle in the actual sequence. Sensitivity so defined is strictly stronger than dependence, it is destroyed by any route through a state that does not separate contents calling for different movements, and it is preserved by any route that does. The climber, the party guest, the stampeding pigs and the reckless nephew fail it; Frankfurt's Jones, the prosthetic patient, the ideomotor reacher and the skilled driver pass it. The right way for a reason to cause a movement is the way that makes the movement sensitive to counterfactual variations in the reason.

\begin{thebibliography}{99}\small
\bibitem[Aguilar(2012)]{Aguilar2012} Aguilar, J.~H. (2012). Basic causal deviance, action repertoires, and reliability. \emph{Philosophical Issues}, 22, 1--19.
\bibitem[Anscombe(1957)]{Anscombe1957} Anscombe, G.~E.~M. (1957). \emph{Intention}. Basil Blackwell.
\bibitem[Bishop(1989)]{Bishop1989} Bishop, J. (1989). \emph{Natural Agency: An Essay on the Causal Theory of Action}. Cambridge University Press.
\bibitem[Brand(1984)]{Brand1984} Brand, M. (1984). \emph{Intending and Acting: Toward a Naturalized Action Theory}. MIT Press.
\bibitem[Chisholm(1966)]{Chisholm1966} Chisholm, R.~M. (1966). Freedom and action. In K.~Lehrer (ed.), \emph{Freedom and Determinism}. Random House.
\bibitem[Davidson(1963)]{Davidson1963} Davidson, D. (1963). Actions, reasons, and causes. \emph{The Journal of Philosophy}, 60(23), 685--700.
\bibitem[Davidson(1973)]{Davidson1973} Davidson, D. (1973). Freedom to act. In T.~Honderich (ed.), \emph{Essays on Freedom of Action}. Routledge \& Kegan Paul.
\bibitem[Frankfurt(1969)]{Frankfurt1969} Frankfurt, H.~G. (1969). Alternate possibilities and moral responsibility. \emph{The Journal of Philosophy}, 66(23), 829--839.
\bibitem[Frankfurt(1978)]{Frankfurt1978} Frankfurt, H.~G. (1978). The problem of action. \emph{American Philosophical Quarterly}, 15(2), 157--162.
\bibitem[Goldman(1970)]{Goldman1970} Goldman, A.~I. (1970). \emph{A Theory of Human Action}. Prentice-Hall.
\bibitem[Hursthouse(1991)]{Hursthouse1991} Hursthouse, R. (1991). Arational actions. \emph{The Journal of Philosophy}, 88(2), 57--68.
\bibitem[Hyman(2015)]{Hyman2015} Hyman, J. (2015). \emph{Action, Knowledge, and Will}. Oxford University Press.
\bibitem[James(1890)]{James1890} James, W. (1890). \emph{The Principles of Psychology}. Henry Holt.
\bibitem[Melden(1961)]{Melden1961} Melden, A.~I. (1961). \emph{Free Action}. Routledge \& Kegan Paul.
\bibitem[Mele(1992)]{Mele1992} Mele, A.~R. (1992). \emph{Springs of Action: Understanding Intentional Behavior}. Oxford University Press.
\bibitem[Peacocke(1979)]{Peacocke1979} Peacocke, C. (1979). \emph{Holistic Explanation: Action, Space, Interpretation}. Clarendon Press.
\bibitem[Schlosser(2007)]{Schlosser2007} Schlosser, M.~E. (2007). Basic deviance reconsidered. \emph{Analysis}, 67(3), 186--194.
\bibitem[Searle(1983)]{Searle1983} Searle, J.~R. (1983). \emph{Intentionality: An Essay in the Philosophy of Mind}. Cambridge University Press.
\bibitem[Sehon(2005)]{Sehon2005} Sehon, S. (2005). \emph{Teleological Realism: Mind, Agency, and Explanation}. MIT Press.
\bibitem[Shepherd(2021)]{Shepherd2021} Shepherd, J. (2021). \emph{The Shape of Agency: Control, Action, Skill, Knowledge}. Oxford University Press.
\bibitem[Steward(2012)]{Steward2012} Steward, H. (2012). \emph{A Metaphysics for Freedom}. Oxford University Press.
\bibitem[Thompson(2008)]{Thompson2008} Thompson, M. (2008). \emph{Life and Action: Elementary Structures of Practice and Practical Thought}. Harvard University Press.
\bibitem[Velleman(1992)]{Velleman1992} Velleman, J.~D. (1992). What happens when someone acts? \emph{Mind}, 101(403), 461--481.
\bibitem[Wilson(1989)]{Wilson1989} Wilson, G.~M. (1989). \emph{The Intentionality of Human Action}. Stanford University Press.
\bibitem[Wittgenstein(1953)]{Wittgenstein1953} Wittgenstein, L. (1953). \emph{Philosophical Investigations}. Trans. G.~E.~M. Anscombe. Basil Blackwell.
\bibitem[Woodward(2003)]{Woodward2003} Woodward, J. (2003). \emph{Making Things Happen: A Theory of Causal Explanation}. Oxford University Press.
\end{thebibliography}

\end{document}
